\documentclass{situationist-preprint}

\title{Popper's Criterion Applied to Itself}
\msid{SCI-07-b}
\msfamily{SCI-07}
\msdiscipline{Philosophy of Science}
\mstype{Companion manuscript}
\mskeywords{Popper; falsifiability; self-refutation; methodological rules; Lakatos}

\begin{document}
\maketitle

\begin{abstract}
SCI-07-a argues that the falsifiability criterion is not self-refuting, because it is a methodological proposal outside the scope it governs. We show that this was Popper's own understanding of the criterion, and we examine the most serious attempt to apply it to itself, Lakatos's proposal to test methodologies against the history of science. That proposal concerns a different object: not the criterion, but the empirical hypothesis that scientists have followed it. The hypothesis is falsifiable and may be false, while the criterion remains unfalsifiable and unrefuted.
\end{abstract}

\section{Popper's view of his criterion}\label{s:popper}

\citet[§4]{Popper1959} presented his criterion of demarcation as a proposal for an agreement or convention, to be judged by its consequences and not by its truth, and he described the rules of empirical method as conventions, comparable to the rules of a game \citep[§§10--11]{Popper1959}. A proposal is not a statement put forward as belonging to empirical science. This is Assumption~2.3 of SCI-07-a, and Theorem~3.2 of SCI-07-a then shows that the falsifiability of the criterion makes no difference to whether its norm is satisfied. The self-application charge therefore does not touch Popper's criterion as Popper understood it.

\section{The historiographic counterpart}\label{s:lakatos}

\citet{Lakatos1971} proposed that a methodology be appraised by the rational reconstruction of history it yields: a methodology is refuted, in a sense, if it condemns as irrational the episodes the scientific elite judges to be the best. This looks like an application of falsifiability to falsificationism. It is not.

\begin{proposition}\label{p:counterpart}
The hypothesis that the history of science conforms to the falsifiability criterion is an empirical statement within the scope of the criterion; the criterion is not. The refutation of the hypothesis is not a refutation of the criterion.
\end{proposition}

\begin{proof}
The hypothesis is a claim about what scientists have done, and observations of the history of science could tell against it, so it is falsifiable and within the scope of the criterion. The criterion is a rule, and a rule is not refuted by a record of its being broken. If scientists have often retained theories in the face of refutation, as \citet{Kuhn1970} argued, the hypothesis is false; the criterion remains a proposal about what scientists should do, whose merits are to be judged by the consequences of adopting it.
\end{proof}

Lakatos's appraisal of methodologies is legitimate, but it bears on the consequences by which a proposal is judged, not on its truth. A proposal that would condemn the whole of successful science would have bad consequences and should be revised; that is an argument against adopting it, not a demonstration that it refutes itself.

\begin{objection}[Verification]
The verificationist could make the same move and call the verification principle a proposal.
\end{objection}

\begin{reply}
The verification principle governs all meaningful statements, and if it is put forward as meaningful it falls within its own scope, where Proposition~3.1 of SCI-07-a applies. The falsifiability criterion governs only the statements of empirical science, and it has a place to stand outside them. The difference is in the scopes, and so the same move does not succeed for both.
\end{reply}

\section{Conclusion}\label{s:conclusion}

Applied to itself, Popper's criterion yields no contradiction, because it is a proposal outside its scope, as Popper said. What can be tested is the historical hypothesis that science conforms to it, and the fate of that hypothesis is a consideration in favor of or against adopting the proposal, not a refutation of it.

\begin{thebibliography}{9}
\bibitem[Kuhn(1970)]{Kuhn1970} Kuhn, T.~S. (1970). \emph{The Structure of Scientific Revolutions}. 2nd edition. University of Chicago Press.
\bibitem[Lakatos(1971)]{Lakatos1971} Lakatos, I. (1971). History of science and its rational reconstructions. In R.~C. Buck and R.~S. Cohen (eds.), \emph{PSA 1970}, Boston Studies in the Philosophy of Science, vol.~8, 91--136. Reidel.
\bibitem[Popper(1959)]{Popper1959} Popper, K.~R. (1959). \emph{The Logic of Scientific Discovery}. Hutchinson.
\end{thebibliography}

\end{document}
